\documentclass[11pt]{article}
\usepackage[margin=1in]{geometry}
\usepackage[T1]{fontenc}
\usepackage{lmodern}
\usepackage{microtype}
\usepackage{natbib}
\usepackage{float}
\usepackage{amssymb}
\usepackage{../shared/smcbvs}
\bibliographystyle{abbrvnat}
\setcitestyle{authoryear,round}
\linespread{1.02}
\setlength{\emergencystretch}{2em}
% assumption labels of "What we prove", printed without links (separate document)
\newcommand{\lab}[1]{\textup{(\asmdisp{#1})}}
\newcommand{\SL}{Song and Liang}
% paragraph heads: name, labels, verdict
\newcommand{\ahead}[3]{\par\medskip\noindent\textbf{#1} #2. \emph{#3}\par}
\newcommand{\ev}[1]{\noindent\emph{Evidence.} #1\par}
\newcommand{\ro}[1]{\noindent\emph{Rules out.} #1\par}
\newcommand{\ds}[1]{\noindent\emph{Our designs.} #1\par}
\newcommand{\tu}[1]{\noindent\emph{A typical user.} #1\par}
\newcommand{\yes}{$\bullet$}\newcommand{\no}{$\circ$}\newcommand{\unc}{?}\newcommand{\na}{--}

\begin{document}
\begin{center}
{\Large\bfseries Are the assumptions reasonable?}\par\smallskip
{\small SMC-BVS project. Draft, \today}
\end{center}

% =====================================================================================================
\section{The answer}\label{sec:a-answer}

Our theorems use the same kinds of assumptions as \citet{SongLiang2015}: a Gaussian linear model, exact sparsity, a
bound on the dimension, a law for the predictors, eigenvalue floors on small sets of columns, and nonzero coefficients
that are not too small. On the designs we study, the assumptions of our posterior theorems hold, strong factors
included. Where the two papers differ, the direction depends on the regime --- their model-size cap, whether predictors
are centred, and which of their theorems is the comparator --- and is given assumption by assumption in the table.

For the method as coded, the theorem that the run reports the true model holds on our independent and locally
correlated designs, and not on pervasive common factors when $p$ is much larger than $\sqrt n$. The theorem that the
printed estimates are consistent needs one more condition, \emph{no leakage}, which only our independent designs
satisfy. On every design the posterior theorems hold and the evidence estimate is unbiased. Conclusions are compared
with Song and Liang theorem by theorem (Section~\ref{sec:a-setting}), not as a whole.

% =====================================================================================================
\section{Theorems, designs, and what ``inside'' means}\label{sec:a-setting}

\paragraph{The theorems.} Numbers refer to the companion section ``What we prove''. The \emph{posterior theorems} are
the \emph{selection theorems} (the posterior finds the true model; Theorems~3.1, 3.4 and~3.9) and the
\emph{distribution theorem} (it finds the true conditional law of the response; Theorem~4.2). In primitive terms,
Theorem~3.9 needs only a Gaussian linear model with i.i.d.\ Gaussian rows; a correlation matrix whose smallest
eigenvalue is bounded below (its largest is free); a fixed number of nonzero coefficients that do not change with $n$;
and $p\le n^A$. The \emph{reported-model theorem} (Theorem~5.7) says that the method, as coded, reports the true model;
it adds three conditions that come from how the method searches (Section~\ref{sec:a-rm}). The \emph{printed-estimates
theorem} (Theorem~5.19) says that the printed inclusion probabilities, median model and estimate of the evidence (the
normalising constant of the posterior) are consistent on the event that the method uses the sampler and its
calibration succeeds; it asserts no lower bound on the probability of that event, and exact enumeration is outside it.
The evidence estimate is unbiased on every design, for each of the four independent sampler runs that the method
pools and for the pooled estimate (status of the proofs: Section~\ref{sec:a-math}).

\paragraph{How the conclusions compare.} Against their Theorem~3 (p.~956), our selection theorems add specificity and
exact selection at a polynomial rate, while theirs gives sensitivity with an exponentially small error; neither implies
the other. Against their Theorem~S (supp.\ p.~13), the statements concern different posteriors; within one posterior,
mass on the true model with a rate is a stronger form of statement than a statement about the posterior mode. The
distribution theorem uses the same distance as their Theorems~1--2, with the same exponential form of the tail for
levels $M\le\sqrt n/\log n$; it reaches a radius of order $\rn$, which their conditions do not reach at polynomial $p$,
and their theorems reach larger radii and exponential $p$, which ours does not. For the reported model, their
statements are limits in the number of iterations at fixed data, for every data set; ours hold at the coded budget
under design conditions; neither implies the other.

\paragraph{The designs.} All our designs have i.i.d.\ Gaussian rows, noise variance $1$ and fixed nonzero coefficients.
They fall into three classes.
\begin{itemize}[leftmargin=1.5em,itemsep=1pt]
  \item \emph{Independent:} S1 of the 2,400-fit study, and IND of the $10^6$ plan (planned runs with $p$ up to $10^6$).
  \item \emph{Locally correlated:} S2 (AR(1), correlation $0.5$), S3, S4 (blocks of five at $0.7$, $0.85$), S6 (blocks of
    five at $0.4$--$0.9$), AR (correlation $0.6$, the primary design of the $10^6$ plan) and BLK (blocks of $20$ at
    $0.5$).
  \item \emph{Pervasive factor:} S5 (blocks of five at $0.7$ plus a global factor, $0.3$ between all columns); the
    strong-factor study (one factor at $0.3$, $0.5$ or $0.7$; two factors; and its own version of S5; true effects of
    aligned signs throughout), with $n$ from $50$ to $150$ and $p$ up to $10n$; and FAC (one factor at $0.3$), a
    contingent design of the $10^6$ plan, run at $p=10^4$ only.
\end{itemize}
The 2,400-fit study has $n=200$, $p\in\{500,1000\}$ and $14$ true predictors, one in each group of five; the $10^6$ plan
has $5$ true predictors.

\paragraph{What ``inside'' means.} A design is \emph{inside} an assumption when the design, continued as a sequence in
$n$ with its structure and the number of true predictors fixed, satisfies it. We judge the qualitative conditions (a
floor bounded below, nonzero coefficients that do not change with $n$, a leakage set of size $n^b$); the explicit
constants of the proofs are conservative and are not part of the judgement. Page numbers are journal pages of their paper; ``supp.\ p.'' is a page of its online supplement; $\bar r_n$ is
their model-size cap.

\begin{table}[tp]
\centering\footnotesize
\setlength{\tabcolsep}{3.5pt}
\renewcommand{\arraystretch}{1.1}
\caption{The assumptions, \SL{} (2015), and our designs, by class: independent, locally correlated, pervasive factor.
\yes{} inside, \no{} outside, \unc{} unclear, \na{} not applicable. A comparison in one word appears only where it holds
in every regime of both papers.}\label{tab:a-summary}
\begin{tabular}{@{}>{\raggedright\arraybackslash}p{0.19\linewidth}>{\raggedright\arraybackslash}p{0.17\linewidth}>{\raggedright\arraybackslash}p{0.41\linewidth}ccc@{}}
\toprule
Assumption & In words & \SL: counterpart, and comparison by regime & Indep. & Local & Factor \\
\midrule
\multicolumn{6}{@{}l}{\textbf{All theorems}} \\
Gaussian linear model \lab{GL} & Gaussian noise, fixed variance & Model (1), density (2), p.~950. \emph{The same, in every
regime.} & \yes & \yes & \yes \\
Dimension \lab{DIM} & $p\le n^A$ & Cor.~1, p.~954; supp.\ p.~13 and Cor.~S$_1$, supp.\ p.~16; $P_n\succ n^\theta$,
p.~950. For large $p$ ours admits fewer dimensions; ours admits bounded $p$, theirs does not. & \yes & \yes & \yes \\
Fixed true model size \lab{SP} & $\sz$ fixed & Strict sparsity (4), p.~950. \emph{The same} for their Theorems 2, 3 and
S; their Theorem~1 also admits growing supports under (3). & \yes & \yes & \yes \\
\midrule
\multicolumn{6}{@{}l}{\textbf{Posterior theorems}} \\
Signal size & bounded; nonzero, not changing with $n$ (beta-min in Thm~3.1 only) & No lower bound stated; (3), (8),
pp.~950, 952; proof rates, supp.\ pp.~6, 15. Depends on whether their $\bstar$ may vary with $n$. & \yes & \yes & \yes \\
Eigenvalue floor \lab{EF}, \lab{EFS} & fixed floor on small sets & Cond.~3, p.~955; (22)--(23), supp.\ p.~13. Where their
sets are no larger than ours, ours implies theirs; otherwise neither implies the other. & \yes & \yes & \yes \\
Near-exact collinearity \lab{NC} & floor $c_\psi/n$ on $O(\log n)$ sets & None in their Theorems~1--2; Cond.~3, p.~955;
(22), supp.\ p.~13. Against the last two, the direction depends on $\bar r_n$. & \yes & \yes & \yes \\
Law of the rows \lab{GR} & Gaussian or sub-Gaussian & $\abs{x_i}\le1$, p.~952; Cond.~1, p.~953. On centred, unit-variance
predictors ours implies Cond.~1; otherwise neither implies the other. & \yes & \yes & \yes \\
\midrule
\multicolumn{6}{@{}l}{\textbf{Reported-model theorem, in addition}} \\
Visibility \lab{V} & no true predictor cancelled by $\le15$ others & Cond.~5, p.~958. \emph{Ours implies theirs, in every
regime.} & \yes & \yes & \unc$^{a}$ \\
Polynomial leakage set, 5.6(a) & at most $n^b$ imitators of the signal & \emph{No hypothesis of this kind}; informal
remark (a), p.~958. No ordering. & \yes & \yes & \no \\
Floor and local regressions, 5.6(b),(c) & each regression in the set uses few members & Floor: as the eigenvalue floor,
on other sets. Local regressions: \emph{no counterpart}. & \yes & \yes & \na \\
\midrule
\multicolumn{6}{@{}l}{\textbf{Printed-estimates theorem, in addition}} \\
No leakage \lab{NL} & null predictors nearly uncorrelated with true ones & \emph{No counterpart}, in every regime. & \yes
& \no & \no \\
\bottomrule
\end{tabular}
\par\smallskip\raggedright\footnotesize
$^{a}$ Inside for every design of the strong-factor study (true effects of aligned signs); unclear for S5 and FAC.
\end{table}

% =====================================================================================================
\section{The assumptions one by one}\label{sec:a-each}

Every design is inside each assumption of Sections~\ref{sec:a-all} and~\ref{sec:a-post}; ``Our designs'' lines
appear only in Sections~\ref{sec:a-rm} and~\ref{sec:a-pe}. Each paragraph opens with the comparison with \SL.

\subsection{Assumptions of all theorems}\label{sec:a-all}

\ahead{Gaussian linear model}{\lab{GL}}{The same as \SL{} in every regime (their model (1) and density (2), p.~950).}
\ev{Their density is $f^*(y,x)=\phi(y;x^\top\bstar,\sigs)\,\nu_x(dx)$; ours is the same model with an intercept.}
\ro{Heavy-tailed or heteroscedastic errors; no theorem of ours or of theirs covers them.}

\ahead{Dimension $p\le n^A$}{\lab{DIM}}{For large $p$, ours admits fewer dimensions than \SL: far fewer than their
Corollary~1 ($P_n<\exp(Cn^\alpha)$, $\alpha\in(0,1)$, p.~954), and slightly fewer than their Theorem~S (supp.\ p.~13);
for small $p$, ours admits bounded $p$, and theirs does not ($P_n\succ n^\theta$, p.~950).}
\ev{Their Theorem~S needs $\log(\sqrt n/r_n)\succ\sqrt{\log P_n}$ (supp.\ p.~13, condition~2), hence
$\log P_n\prec(\log n)^2/4$; their Corollary~S$_1$ states $P_n<n^{\kappa\log n}$ with $\kappa\in(0,\tfrac14)$ (supp.\
p.~16). Exponential $p$ is not covered. Our largest exponent is the $10^6$ plan at
$n=200$: $A=\log10^6/\log200\approx2.6$.}
\ro{Numbers of predictors that grow faster than every power of $n$, such as $\log p\asymp n^\alpha$.}

\ahead{Fixed true model size}{\lab{SP}}{The same as \SL's strict sparsity (4), p.~950, which their Theorems~2, 3 and~S
use; their Theorem~1 alone also allows (3), p.~950, under which the number of nonzero coefficients may grow, a case
ours does not cover.}
\ev{It cannot be dropped: on every design with $\sz\le\min(K-1,p/2)$, $\pis{\Sz}\le\{1+(\sz+1)/(3\sqrt n)\}^{-1}$
(Theorem~3.13(a)), so a posterior miss of order $\rn$ is possible only if $\sz=O(\sqrt{\log n})$.}
\ro{True models whose size grows with $n$.}

\subsection{Assumptions of the posterior theorems}\label{sec:a-post}

\ahead{Signal size}{\lab{BC}, \lab{FSb}, \lab{FSt} and Assumption~4.1(e), the bounds; \lab{BM}, the beta-min}{\SL{}
state no lower bound on the nonzero coefficients. If their $\bstar$ is one fixed sequence, the natural reading of
p.~950, our Theorems~3.9, 5.7 and~5.19 ask the same, and Theorems~3.1 and~4.2, whose coefficients may vary with $n$,
admit more; if their $\bstar$ may vary with $n$, their proof steps need rates (supp.\ pp.~6, 15), and our fixed
signals ask more than those steps use.}
\ev{Our explicit beta-min \lab{BM}, at the usual order $\sigs\rn$,
appears only in the fixed-design Theorem~3.1; it has the same form as their proof rates, with a fixed multiple of
$\log n$ where theirs has a diverging multiple of $\bar r_n\log P_n$ (supp.\ p.~15). Their Theorem~1 bounds the
coefficients above through (3) and (8) (pp.~950, 952). Without any lower bound, specificity can fail
(Remark~3.2(ii)).}
\ro{True coefficients below a constant times $\rn$ (Theorem~3.1 only), and signals that grow without bound.}

\ahead{Eigenvalue floor on small sets}{\lab{EF}, \lab{EFS}, Assumptions~3.3 and~4.5(a)}{Where \SL's sets are no
larger than ours ($\abs t+\bar r_n\le2(A+1)\sz+2$ against Theorem~3.1; $\bar r_n+1\le5(A+1)(\sz+1)\log n$ against
Theorem~3.9), our fixed floor implies their Condition~3 (p.~955) and their (22)--(23) (supp.\ p.~13), and also asks a
floor on sets where they ask nothing; where their sets are larger ($\bar r_n\succ\log n$), they ask a vanishing level on
sets where we ask nothing, and neither condition implies the other.}
\ev{Condition~3 asks a \emph{vanishing} level $\delta_n\succ\epsilon_n^2$ on $\bar r_n+1$ columns that contain a true
predictor. We ask
a \emph{fixed} level on sets of $(2A+2)\sz+2$ predictors (fixed design) or $O(\log n)$ predictors (Gaussian rows). In
every regime both conditions follow from $\lmin(\Sig)\ge c$. No upper eigenvalue bound is used (the upper half of
their (22) is no restriction for standardised columns), so pervasive factors are allowed: the smallest eigenvalue of
$\Sig$ is $1$ (S1, IND), $1/3$ (S2), $0.3$ (S3, S5), $0.15$ (S4), $0.1$ (S6), $1-\rho$ (one factor at $\rho$), $0.40$ (two
factors), $0.25$ (AR), $0.5$ (BLK) and $0.7$ (FAC).}
\ro{Exact or asymptotic duplicates of a true predictor, and correlations tending to one with a true predictor.}

\ahead{No near-exact collinearity on sets of logarithmic size}{\lab{NC}, Assumption~4.1(d)}{\SL's Theorems~1--2 have
no condition of this kind, in every regime; against their Condition~3 (p.~955) and (22) (supp.\ p.~13), theirs gives
ours on the sets they cover if $\bar r_n\ge3(A+1)(\sz+1)\log n$, and ours constrains sets that theirs does not if
$\bar r_n\prec\log n$, so that neither implies the other.}
\ev{Their eigenvalue conditions (11) and (14) (p.~952) are on the prior covariance. With their own ridge choice, a ridge
$c_n\propto1/n$ (p.~953), no design condition is left, at the price of a slower rate (their (10)). Our floor, $c_\psi/n$
on the correlation scale, is of the same order as that ridge, on $O(\log n)$ sets only, and gives the distribution
theorem its contraction radius $\sqrt{M\log n/n}$, for a fixed level $M$. Without it the distribution theorem still
holds, with an extra factor $\sqrt{\log(\praw\vee n)}$ in the radius, $\praw$ being the number of raw predictors
(Proposition~4.3(ii)): the assumption costs one log factor.}
\ro{Duplicated columns (null duplicates included) and exact linear dependence, as in one-hot or compositional coding.}
\tu{Inside after removing exact duplicates and exact linear dependences.}

\ahead{Law of the rows}{\lab{GR}, Assumption~4.1(c)}{On centred, unit-variance predictors, Gaussian rows and
Assumption~4.1(c) each imply \SL's Condition~1 (p.~953); on raw predictors, and against their $\abs{x_i}\le1$
(Theorem~1, p.~952), neither implies the other.}
\ev{For standardised predictors, Assumption~4.1(c) implies their Condition~1 for every $r$, with $\delta=\sqrt{3/8}/L$
and $M=\log2$, $L$ being the sub-Gaussian constant of Assumption~4.1(c) (a two-line computation, independently
checked). On raw predictors, Gaussian rows with mean $n$
satisfy our row law and violate Condition~1, which is a raw exponential-moment bound. Our generators draw centred
Gaussian rows with unit variances, which satisfy Condition~1 itself.}
\ro{Non-Gaussian rows for Theorems~3.9, 5.7 and 5.19; heavy-tailed predictors for the distribution theorem.}
\tu{Bounded data such as genotypes are outside the Gaussian-row theorems as stated, and inside \SL's Condition~1; they are
covered by the fixed-design Theorem~3.1 when the observed design meets its conditions.}

\subsection{Further assumptions of the reported-model theorem}\label{sec:a-rm}

\ahead{Visibility}{\lab{V}}{Ours implies \SL's Condition~5 (p.~958: $\abs{\operatorname{cov}(x_t,y)}>0$ for every true
predictor) in every regime: theirs is our condition with no conditioning and no margin.}
\ev{Visibility asks $\abs{\Cov(W_j,W^\top\tstar\mid W_U)}\ge a>0$ for every true predictor $j$, with $W$ the standardised
predictors and $U$ any set of at most $15$ other columns: the start law (the distribution the sampler starts from) must
see every true predictor, as their first stage must, and it conditions on up to $15$ columns because the start law works on groups of at most $16$ predictors. For block and AR designs with at most one true
predictor per correlation block, this covariance is $\theta^*_j\Var(W_j\mid W_U)$ plus geometrically small terms, at
least $\abs{\theta^*_j}\lmin(\Sig)/2$. When the true effects have aligned signs and every conditional covariance
between columns outside $U$ is nonnegative, as in each design of the strong-factor study, it is at least
$\theta^*_j\lmin(\Sig)$.}
\ro{Masking: a true predictor whose signal some $15$ other predictors cancel.}
\ds{Inside: the independent and locally correlated designs, and every design of the strong-factor study. Unclear: S5
(mixed signs; a search found a near-cancellation) and FAC (a search found a margin below the aligned-sign bound, so its
signs are mixed).}
\tu{A numerical check of the conditions' population quantities, not yet complete, found visibility failing for a
pervasive factor with effects of mixed signs and for a riboflavin-like design; such a user may be outside this theorem
and still inside the posterior theorems.}

\ahead{Polynomial leakage set}{Assumption~5.6(a)}{\SL{} have no hypothesis of this kind; the nearest text is their
informal remark (a), p.~958, which is not a condition of any of their theorems. No ordering is claimed.}
\ev{They write: ``if there are $p_n$ predictors correlated with the true predictor $X_t$, $p_n\prec P_n$, then there are
at most $cp_n$ predictors to be selected'' (p.~958). Their results on their sampler's output carry no design condition of
this kind (supp.\ pp.~9--10, 12--13, 16--17), so there is no formal counterpart to compare levels with. We state the requirement as a hypothesis with a rate: the null predictors
whose conditional covariance with a true predictor exceeds a level $\varpi_n=o(\rn)$ must lie in a set of at most $n^b$
columns, the leakage set $Q$, with $b<\tfrac12$.}
\ro{Pervasive factors with $p\gg\sqrt n$: with no conditioning, every null predictor then has a fixed covariance with every
true one, so the leakage set is all of $\set{p}$, and $\abs{Q}\le n^b$ would force $p\le n^b$.}
\ds{\textbf{Inside:} S1 (no leakage set beyond the true model); S2 and AR (a neighbourhood of size $C\log n$ around the true
predictors); S3, S4, S6 and BLK (the true predictors' blocks). \textbf{Outside:} S5, the whole strong-factor study, and
FAC.}

\ahead{Floor and local regressions inside the leakage set}{Assumption~5.6(b),(c)}{The floor on $Q\cup\{l\}$ is an
eigenvalue floor of the kind of \SL's Condition~3 (p.~955) and (22) (supp.\ p.~13), on a different family of sets, and
compares with theirs as the eigenvalue floor does; the local regressions have no counterpart in their paper (they serve
the search, which they do not analyse).}
\ev{Both follow from the global floor \lab{EFS} plus a locality property of $\Sig$: for block, banded and AR structures,
take as local set the block-mates or the nearest leakage-set neighbours. A fixed-design version needs nothing of the kind
when the leakage set has at most $(1-\delta)\log n$ extra predictors, for a fixed $\delta\in(0,1)$ (Theorem~5.12, case
(d1)).}
\ro{Leakage sets collinear with another predictor, or inside which regressions depend on many members.}
\ds{As for the polynomial leakage set; not applicable where that fails.}

\subsection{Further assumption of the printed-estimates theorem}\label{sec:a-pe}

\ahead{No leakage}{\lab{NL}}{\SL{} have no condition of this kind, in every regime: none of their theorems concerns a
start law. Only our independent designs satisfy it.}
\ev{With no conditioning, \lab{NL} asks $\abs{\Cov(W_l,W_j)}\le C_Ln^{-1/2}$ for every null predictor $l$ and true
predictor $j$. In AR(1) with correlation $0.5$ the neighbour of a true predictor has covariance $0.5$; in a block design
every block-mate of a true predictor has covariance at least $0.4$. For such designs the start law concentrates on the
models between the true model and the leakage set, not on the true model alone. On AR and BLK the present proof route
cannot be extended as $n\to\infty$, because the start law's mass on the true model then decays faster than any power of
$n$; at the sizes of the $10^6$ plan that result does not apply, and nothing recorded bounds this mass there in either
direction. For S3, S4 and S6, a theorem under a weaker, group-relative form of \lab{NL} is a candidate; it is not
proved.}
\ro{Any null predictor correlated with a true one at a fixed level.}
\ds{\textbf{Inside:} the independent designs only. \textbf{Outside:} all locally correlated and factor designs, including AR, the
primary design of the $10^6$ plan.}

% =====================================================================================================
\section{Assumptions on mathematical objects}\label{sec:a-math}

Premises about the prior, distributions, distances, Markov kernels and the code must be proved, not argued.
\begin{itemize}[leftmargin=1.5em,itemsep=2pt]
  \item \emph{Posterior theorems.} The prior facts, the posterior-summary lemma (Lemma~2.1) and the Hellinger facts
    are proved and independently checked, and the output description used by Corollary~4.7 is checked against the code.
    Two lemmas behind the expectation bound of Theorem~3.1 await their independent check, and one link in the record of
    Corollary~4.7 is to be corrected.
  \item \emph{Theorems on the coded method.} The probabilistic events of the reported-model proof are proved and
    independently checked. Thirteen premises are named and queued. Ten are code readings or short mathematical checks,
    for example that the run guard stops runs without changing them, that the start sampler is exact, and that the coded
    moves are the stated kernels. Of the other three, one is used only by Remark~5.21 and needs a new key lemma; two
    are used only by an older proof that the core theorems no longer use. The unbiasedness
    of the evidence estimate also rests on two inputs not yet checked under the current protocol: the identities of the
    frozen final run, and the two-line argument for the pooled form.
  \item \emph{Idealisations, stated as conventions:} exact arithmetic and ideal random numbers.
\end{itemize}

\renewcommand{\bibsection}{\par\medskip\noindent\textbf{Reference}\par}
{\small\bibliography{../shared/refs}}
\end{document}
